\section*{Chapter \ref{ch:b2:conjugate-additions} --- Conjugated Carbonyls: Michael and Robinson}

\begin{solution}{exo:b2:conjugate-additions:1}
\begin{itemize}
  \item \ce{CH3MgBr}: mainly 1,2 (hard, irreversible), 1-methylcyclohex-2-en-1-ol.
  \item \ce{(CH3)2CuLi}: 1,4, 3-methylcyclohexanone.
  \item Thiophenol with base: 1,4 (the thiolate is soft), 3-(phenylsulfanyl)cyclohexanone.
  \item \ce{LiAlH4}: 1,2, cyclohex-2-en-1-ol.
\end{itemize}
\end{solution}

\begin{solution}{exo:b2:conjugate-additions:2}
Donors: pentane-2,4-dione, nitromethane, diethyl propanedioate (acidic \ce{C-H}). Acceptors:
propenenitrile, methyl propenoate, but-3-en-2-one (\ce{C=C} conjugated with an acceptor group).
\end{solution}

\begin{solution}{exo:b2:conjugate-additions:3}
\ce{C4H9Br + 2Li -> C4H9Li + LiBr}, then \ce{2C4H9Li + CuI -> (C4H9)2CuLi + LiI}, in ether at low
temperature under inert gas.
\end{solution}

\begin{solution}{exo:b2:conjugate-additions:4}
Diethyl (3-oxobutyl)propanedioate, \ce{(EtOOC)2CHCH2CH2COCH3}. Hydrolysis gives the substituted
propanedioic acid, which loses \ce{CO2} on heating (\cref{prop:b2:enolates-aldol:decarboxylation}):
5-oxohexanoic acid, \ce{HOOCCH2CH2CH2COCH3}.
\end{solution}

\begin{solution}{exo:b2:conjugate-additions:5}
Lithium dimethylcuprate in ether at low temperature, then aqueous work-up: \omterm{def:b2:conjugate-additions:addition-modes}{conjugate addition}.
Methylmagnesium bromide, harder, adds mostly 1,2 and gives 1-methylcyclohex-2-en-1-ol as the main
product.
\end{solution}

\begin{solution}{exo:b2:conjugate-additions:6}
\ce{Et3N} removes a little \ce{EtSH} to give \ce{EtS-}; the thiolate adds to the $\beta$ carbon, the
\omterm{def:b2:enolates-aldol:enolate}{enolate} formed takes the proton of another \ce{EtSH}, which regenerates \ce{EtS-}. Product:
4-(ethylsulfanyl)butan-2-one. Sulfur is large and polarisable, its lone pair high in energy: a \omterm{def:b2:frontier-orbitals:hard-soft}{soft
nucleophile}, under \omterm{def:b2:frontier-orbitals:control}{orbital control}, so it attacks where the \omterm{def:b2:frontier-orbitals:frontier}{LUMO} is largest, and the addition is
reversible, which also favours the more stable 1,4-adduct.
\end{solution}

\begin{solution}{exo:b2:conjugate-additions:7}
Cyanide attacks the more positive carbonyl carbon faster: the cyanohydrin is the kinetic product. Its
formation is reversible; at higher temperature, given time, cyanide is released and re-adds at the
$\beta$ carbon, giving the adduct that keeps the stronger \ce{C=O} $\pi$ bond, the thermodynamic product.
\end{solution}

\begin{solution}{exo:b2:conjugate-additions:8}
With 2-methylcyclohexane-1,3-dione: the Wieland--Miescher ketone (figure of the \omterm{def:b2:conjugate-additions:robinson}{Robinson annulation}).
With cyclohexanone: 4,4a,5,6,7,8-hexahydronaphthalen-2(3\textit{H})-one, a bicyclic enone whose
\ce{C=C} ends at a ring-fusion carbon, with no angular methyl group.
\end{solution}

\begin{solution}{exo:b2:conjugate-additions:9}
\ce{CH3COCH2CH2CH2COCH3}: C2 and C6 are the carbonyls, five carbons from C2 to C6 counted inclusively,
a 1,5-diketone. Cut C3--C4: donor, the \omterm{def:b2:enolates-aldol:enolate}{enolate} of propanone at C3 (better, ethyl 3-oxobutanoate, whose
ester is removed later by hydrolysis and \omterm{def:b2:enolates-aldol:decarboxylation}{decarboxylation}); acceptor, but-3-en-2-one. Conditions:
ethyl 3-oxobutanoate, but-3-en-2-one, catalytic sodium ethoxide in ethanol; then aqueous acid and
heat.
\end{solution}

\begin{solution}{exo:b2:conjugate-additions:10}
The \omterm{def:b2:enolates-aldol:aldol}{aldol} makes a \ce{C-C} bond but turns a \ce{C=O} $\pi$ bond into a \ce{C-O} $\sigma$ bond; the
balance is small and two molecules become one, so $K$ is modest and the \omterm{def:b2:enolates-aldol:aldol}{retro-aldol} is easy. The
\omterm{def:b2:conjugate-additions:michael}{Michael addition} makes a \ce{C-C} bond at the cost of a weaker \ce{C=C} $\pi$ bond and keeps both
\ce{C=O}: it is clearly favourable. The reverse would need the ketone \omterm{def:b2:enolates-aldol:enolate}{enolate} to expel the stabilised propanedioate anion; the
equilibrium lies so far on the adduct side that under catalytic base the adduct persists.
\end{solution}

\begin{solution}{exo:b2:conjugate-additions:11}
\ce{C(COOEt)2(CH2CH2COOCH3)2}. After the first \omterm{def:b2:conjugate-additions:michael}{Michael addition} the central carbon still has one
hydrogen between two esters, as acidic as before: the base removes it and a second addition follows.
\end{solution}

\begin{solution}{exo:b2:conjugate-additions:12}
Through the more substituted \omterm{def:b2:enolates-aldol:enolate}{enolate} (towards the methyl-bearing carbon): the octalone with the methyl
group on the ring-fusion carbon, 4a-methyl-4,4a,5,6,7,8-hexahydronaphthalen-2(3\textit{H})-one. Through
the less substituted \omterm{def:b2:enolates-aldol:enolate}{enolate} (at the \ce{CH2} side): the methyl group ends on a ring carbon next to the
fusion. The first is favoured under equilibrating conditions (an alkoxide in its alcohol, or the \omterm{def:b2:amines:imine}{enamine}
route with a \omterm{def:b2:amines:class}{secondary amine} and acid), which give the thermodynamic, more substituted \omterm{def:b2:enolates-aldol:enolate}{enolate}
(\cref{prop:b2:enolates-aldol:enolate-control}).
\end{solution}

\begin{solution}{pb:b2:conjugate-additions:1}
\textbf{1.} A cyclohexane ring with \ce{C=O} at C1 and C3 and a methyl group on C2; the hydrogen on C2
is the most acidic.
\textbf{2.} Its removal gives an anion whose charge is delocalised over C2 and both oxygens; for
cyclohexanone only one oxygen shares it, and the acidity is too weak to measure in water.
\textbf{3.} Charge on C2; \ce{C1=C2} with \ce{O^-} on C1; \ce{C2=C3} with \ce{O^-} on C3.
\textbf{4.} Yes: with hydroxide, $K = 10^{14.0 - 5.3} = 10^{8.7}$ (water as the conjugate acid, from
$K_e$); with ethoxide, $10^{15.9-5.3} = 10^{10.6}$. Either is complete.
% ledger: ke:25C, pka:ethanol
\textbf{5.} \ce{C1(OH)=C2(CH3)}: C2 still carries one hydrogen, which is all an \omterm{def:b2:enolates-aldol:tautomerism}{enol} needs; the methyl
group only replaces the second one.
\textbf{6.} Soft: the charge is spread over three atoms and the anion is stabilised. It attacks the
$\beta$ carbon (\cref{prop:b2:conjugate-additions:regio}).
\textbf{7.} The \omterm{def:b2:enolates-aldol:enolate}{enolate}'s C2 bonds to the \ce{CH2} end of but-3-en-2-one; the \omterm{def:b2:enolates-aldol:enolate}{enolate} formed takes a
proton: 2-methyl-2-(3-oxobutyl)cyclohexane-1,3-dione.
\textbf{8.} Each new \omterm{def:b2:enolates-aldol:enolate}{enolate} takes a proton from a molecule of dione (or from the solvent) and so
regenerates the donor \omterm{def:b2:enolates-aldol:enolate}{enolate}: the base is not consumed.
\textbf{9.} The \omterm{def:b2:enolates-aldol:enolate}{enolate} is ambident; with a soft carbon electrophile, under \omterm{def:b2:frontier-orbitals:control}{orbital control}, it reacts
at carbon, and the \ce{C-C} adduct keeps two \ce{C=O} bonds, the stronger combination; an
\ce{O}-adduct, if formed, would revert.
\textbf{10.} One: C2, bonded to C1, C3, the methyl and the chain.
\textbf{11.} The chain carbonyl carbon and the ring carbonyl C1 (or C3): \ce{C(=O)}--\ce{CH2}--\ce{CH2}--C2--C1,
five carbons counted inclusively.
\textbf{12.} To consume all the dione, the more valuable reagent; only slightly, because but-3-en-2-one
is very toxic and polymerises.
\textbf{13.} Ring C4 and C6 (next to C3 and C1), the chain \ce{CH2} next to its carbonyl, and the chain's
terminal \ce{CH3}. C2 has none.
\textbf{14.} The \omterm{def:b2:enolates-aldol:enolate}{enolate} of the terminal \ce{CH3} attacks a ring carbonyl: the ring closed counts that
carbon, the chain carbonyl carbon, two \ce{CH2}, C2 and the attacked carbonyl carbon: six atoms.
\textbf{15.} The internal chain \ce{CH2} \omterm{def:b2:enolates-aldol:enolate}{enolate} would close a four-membered ring, strained. A ring
\omterm{def:b2:enolates-aldol:enolate}{enolate} (C4 or C6) attacking the chain carbonyl closes a six-membered ring too, but bridged: that \omterm{def:b2:enolates-aldol:aldol}{aldol}
cannot dehydrate (the double bond would sit at a bridgehead of a small bicyclic system) and reverts,
while the fused \omterm{def:b2:enolates-aldol:aldol}{aldol} is drained by dehydration.
\textbf{16.} The \omterm{def:b2:enolates-aldol:aldol}{aldol} carries \ce{OH} on the former ring carbonyl carbon, now at the ring fusion; loss
of water with a hydrogen of the former \ce{CH3} carbon (now $\alpha$ to the chain ketone) gives the
\ce{C=C} conjugated with that ketone.
\textbf{17.} The new double bond is conjugated with the carbonyl: an E1cB elimination through the
\omterm{def:b2:enolates-aldol:enolate}{enolate}, favoured by heating (\cref{prop:b2:enolates-aldol:aldol-equilibrium}).
\textbf{18.} A saturated ketone and an $\alpha,\beta$-unsaturated ketone (a cyclohexenone).
\textbf{19.} The ring-fusion carbon carrying the methyl group: four different substituents.
\textbf{20.} No: the dione is achiral, and its two carbonyls are attacked equally (they are mirror images
of each other with respect to the chain); with achiral reagents the product is racemic. Chiral amine
catalysts favour one carbonyl and give one enantiomer in excess (Year 3 volume).
\textbf{21.} Dione \ce{C7H10O2}: \qty{126.155}{g/mol}; \ce{C4H6O}: \qty{70.091}{g/mol}; \ce{C11H14O2}:
\qty{178.231}{g/mol}.
\textbf{22.} \ce{C7H10O2 + C4H6O -> C11H14O2 + H2O}: C 11, H 16, O 3 on each side.
\textbf{23.} $10.0/126.155 = \qty{0.07927}{mol} = \qty{79.3}{mmol}$.
\textbf{24.} $0.07927 \times 178.231 \times 0.70 = \boldsymbol{\approx \qty{9.89}{g}}$ of
Wieland--Miescher ketone.
\end{solution}
